\section{\bt{Design of \system}{\system\ 的设计}}
\label{sec:model}
\bi{The design follows the lifecycle in Figure~\ref{fig:architecture}. We first define the objects in Shared knowledge, then explain how Publish creates them, how Lookup and Maintain sustain their use, and how Compile and run aligns validation with execution. The central distinction is between a definition revision, which specifies a computation, and versioned condition evidence, which determines where that computation can be reused.}
{本节沿图~\ref{fig:architecture} 的生命周期展开：先定义共享知识中的对象，再说明发布组件如何创建这些对象、查找与维护组件如何支持持续复用，以及编译与执行组件如何对齐验证和执行。设计的核心区别在于：定义版本规定计算，带版本的条件证据则决定该计算能够在哪些数据状态上复用。}

\subsection{\bt{Shared representation and compilation}{共享表示与编译}}
\label{sec:shared-object}
\bi{A metric revision is represented as}
{指标的一个版本表示为}
\begin{equation}
\label{eq:metric-revision}
m^r=(B,I,C,E,r).
\end{equation}
\bi{Here, $B$ states the business meaning; $I$ records the structured implementation; $C=C(m^r)$ is the set of executable conditions derived from that implementation; $E$ records the learning task, its verified calculation, dependency versions, and publication checks; and $r$ is the revision number. The index in $m_i^r$ distinguishes different metrics, whereas a change from $m^r$ to $m^{r+1}$ changes one metric's implementation. A feature revision $f^r$ uses the same record structure with additional entity and temporal fields (Section~\ref{sec:features}).}
{其中，$B$ 表达业务含义；$I$ 记录结构化实现；$C=C(m^r)$ 为从实现推导的可执行条件集；$E$ 记录学习题、已核实的计算、依赖版本及发布检查；$r$ 为修订号。$m_i^r$ 的下标区分不同指标，$m^r$ 到 $m^{r+1}$ 则表示同一指标实现的变更。特征版本 $f^r$ 采用同一记录结构，并增加实体和时间字段（第~\ref{sec:features} 节）。}

\bi{The implementation $I$ specifies a fact table $F$, an aggregate operator $\mathrm{AGG}$ and its expression $e$, a business-grain key $K$, a directed join tree $\mathcal J$, definition filters $\varphi$, and a time binding $\tau$. Each join records its columns, direction, inner/left type, and required filters. The tree is rooted at $F$ and follows verified many-to-one edges, covering star and snowflake schemas. The time binding identifies either an event-time column or a fact-to-date-dimension join. The filters encode the definition's selection, such as a row type needed to represent one business event; a consumer's reporting period is supplied separately.}
{实现 $I$ 指定事实表 $F$、聚合算子 $\mathrm{AGG}$ 及其表达式 $e$、业务粒度键 $K$、有向连接树 $\mathcal J$、定义过滤 $\varphi$ 和时间绑定 $\tau$。每条连接记录连接列、方向、内连接或左连接类型，以及所需过滤。连接树以 $F$ 为根，沿经验证的多对一边展开，覆盖星型和雪花型模式。时间绑定指定事件时间列，或事实表到日期维度的连接。过滤编码定义本身的选择，例如表示一个业务事件所需的行类型；使用者的报表期间另行提供。}

\bi{Let $J_I(D_s)$ denote the rows obtained by applying the joins, time binding, and definition filters to database state $D_s$. Given a reporting period $p$, compilation produces a read-only SQL query $q=\operatorname{Compile}(I,p)$ with semantics}
{设 $J_I(D_s)$ 为在数据库状态 $D_s$ 上应用连接、时间绑定与定义过滤后得到的行。给定报表期间 $p$，编译生成只读 SQL 查询 $q=\operatorname{Compile}(I,p)$，其语义为}
\begin{equation}
\label{eq:compiled-metric}
q(D_s)=\mathrm{AGG}_{e}\!\left(\sigma_{\tau\in p}\bigl(J_I(D_s)\bigr)\right).
\end{equation}
\bi{The selection $\sigma_{\tau\in p}$ retains rows whose assigned time belongs to $p$. Period differences and rankings compose these aggregates without relearning $I$. Thus, the two revenue agents in Figure~\ref{fig:architecture} share the same implementation and conditions while changing $p$. Exact query results may be cached separately; answers and question parameters are not part of the reusable definition.}
{选择 $\sigma_{\tau\in p}$ 保留时间归属落在 $p$ 内的行。期间差值与排名可在这些聚合上组合，无需重新学习 $I$。因此，图~\ref{fig:architecture} 中两个营业额智能体共享同一实现与条件，仅改变 $p$。完全相同查询的结果可另行缓存；答案和题目参数不属于可复用定义。}

\bi{Definitions also refer to shared join entries and their revisions. These entries retain validated join directions, cardinalities, and filters, while table descriptions expose schema knowledge discovered during exploration. Check results are separate objects because a condition may be required by several otherwise different definitions. This separation implements both forms of sharing in Figure~\ref{fig:architecture}: the same definition feeds several agents, and the same check feeds several definitions.}
{定义还引用共享连接条目及其修订号。这些条目保留经过验证的连接方向、基数与过滤；表描述提供探索中发现的模式知识。检查结果独立存储，因为同一条件可能被多个不同定义需要。这一分离实现图~\ref{fig:architecture} 中两种共享：同一定义供多个智能体使用，同一检查供多个定义使用。}

\subsection{\bt{Learning and checked publication}{学习与受检查的发布}}
\label{sec:publication}
\bi{The Publish component converts a successful session into a candidate revision through model-based extraction, then admits it through deterministic checks (Figure~\ref{fig:architecture}, steps 1--3). A metric session must state its business definition and have an independently verified answer, supplied by a user, report, or analyst. In the benchmark evaluation, an automated oracle compares the final scalar answer with the reference answer on the same snapshot before extraction. The session also submits query references and arithmetic, so the verified scalar can be traced to the executed computation.}
{发布组件先通过模型提取，将成功会话转为候选版本，再通过确定性检查决定是否接纳（图~\ref{fig:architecture}，第 1--3 步）。指标会话须明确业务定义，并获得用户、报表或分析师对答案的独立核实。基准评估中，自动判题器在提取前，将最终标量答案与同一快照上的参考答案比较。会话还提交查询编号与算术过程，使核实后的标量能够追溯到实际执行的计算。}

\bi{The LLM extractor receives the question, SQL and arithmetic trace, rejected queries, and schema and join metadata. It receives no business rows or result values. Its output is a JSON draft containing the proposed meaning and structured fields. The server parses that draft, derives its conditions and dependencies, and constructs $E$ from the session record. Publication therefore derives the evidence from recorded execution and checks the proposed representation independently.}
{LLM 提取器接收题目、SQL 与算术轨迹、被拒绝的查询，以及模式和连接元数据，不接收业务行或结果值。输出是包含业务含义与结构化字段的 JSON 草稿。服务器解析草稿，推导条件和依赖，并从会话记录构建 $E$。因此，发布从所记录的执行中取得证据，并独立检查所提出的表示。}

\bi{Publication establishes three properties. \emph{Source support} requires an explicit business-definition basis (G1), an independently judged answer, and a verified calculation chain (G2). \emph{Structural admissibility} requires valid tables, columns, expressions, joins, and definition filters with question parameters excluded (G3), uniqueness at the fact's stated grain (G4), and a single read-only example query that passes execution review (G5). For resolved filtered or derived inputs, execution review checks query-local uniqueness using the current schema, without updating base-table evidence.}
{发布检查建立三项性质。\emph{来源支持}要求有明确的业务定义依据（G1），答案经过独立判定且计算链可验证（G2）。\emph{结构合法性}要求表、列、表达式、连接与定义过滤合法，并排除题目参数（G3）；事实表满足所述粒度的唯一性（G4）；示例是通过执行审查的单条只读查询（G5）。对可解析的过滤或派生输入，执行审查依据当前模式在查询作用域内检查唯一性，不据此更新原表证据。}

\bi{\emph{Reproduction} then connects the extracted representation to the successful session. The learning dependencies must remain unchanged and the example SQL must reproduce the verified value (G6). SQL compiled from $I$, instantiated with the learning question's parameters, must independently reproduce that value (G7). Numeric comparisons use the learning task's declared precision. G7 is necessary because consumers execute the compiled fields: retaining a correct example is insufficient if extraction omitted a filter or join. Only passing candidates become available through Lookup; the learning record retains the gates and dependency versions used for admission. Feature publication uses column-level reproduction and temporal audits, as described in Section~\ref{sec:features}.}
{随后，\emph{重现}将提取表示与成功会话联系起来：学习依赖须保持不变，示例 SQL 须重现核实值（G6）；从 $I$ 编译、填入学习题参数的 SQL 也须独立重现该值（G7）。数值比较使用学习任务声明的精度。G7 必不可少，因为使用者执行的是结构化字段的编译结果：若提取遗漏了过滤或连接，保留一个正确示例仍不足以支持复用。只有通过检查的候选才通过查找组件对外提供；学习记录保留接纳时使用的检查与依赖版本。特征发布采用逐列重现和时间审计，见第~\ref{sec:features} 节。}

\subsection{\bt{Conditions and validity}{条件与有效性}}
\label{sec:sufficiency}
\bi{Conditions expose the assumptions needed to preserve the computation after data updates. Even with an unchanged schema, restatements, status-change rows, or duplicate loads may violate the grain; reused dimension keys may make a join fan out; and late or reformatted date keys may move facts to another period or to none~\cite{kimball2013toolkit}. A condition $c$ is a predicate over a dependency set $\operatorname{dep}(c)$ of tables. Its query $Q_c$ returns statistics, and its acceptance rule interprets those statistics. We write $c(D_s)$ when that rule holds on state $D_s$.}
{条件显式表达数据更新后保持计算成立所需的前提。即使模式不变，更正、状态流水或重复加载也可能破坏粒度；复用的维度键可能使连接放大行数；迟到或格式变化的日期键可能使事实归入其他期间或无法归入任何期间~\cite{kimball2013toolkit}。条件 $c$ 是依赖表集合 $\operatorname{dep}(c)$ 上的谓词；查询 $Q_c$ 返回统计量，接纳规则解释这些统计量。规则在状态 $D_s$ 上满足时，记为 $c(D_s)$。}

\bi{\system\ derives four kinds of conditions from $I$. \emph{Filtered key uniqueness}, $\operatorname{KeyUnique}(T,K,\varphi)$, checks uniqueness among the rows of $\sigma_\varphi(T)$ whose key components are non-null. \emph{Join cardinality} bounds the matches of each left row to prevent multiplication. \emph{Date join} checks the fact's date-key mapping, or date coverage when time is a column. \emph{Join completeness} bounds the unmatched fraction for the date join and each inner join by its publication baseline plus 0.1 percentage points. For a snowflake path, the denominator counts fact rows reaching the final edge, so a missing popular dimension member is weighted by the facts it affects.}
{\system\ 从 $I$ 推导四类条件。\emph{过滤后键唯一性} $\operatorname{KeyUnique}(T,K,\varphi)$ 检查 $\sigma_\varphi(T)$ 中各键分量非空的行是否唯一。\emph{连接基数}约束每条左侧行的匹配数，防止行被复制。\emph{日期连接}检查事实日期键的映射，时间直接由列表示时则检查日期覆盖。\emph{连接完整性}要求日期连接和每个内连接的未匹配比例不超过发布时基线加 0.1 个百分点。对于雪花型路径，分母统计到达末端连接的事实行，因此热门维度成员的缺失按其影响的事实数计权。}

\bi{These predicates connect the structured computation to its business meaning. The business definition supplies three domain premises: (B1) selected fact rows have non-null keys $K$ whose values correspond one-to-one to business events; (B2) $e$ evaluates to each event's measurement; and (B3) the time binding and definition filters express the intended business selection. The data checks establish the structural properties needed under these premises.}
{这些谓词将结构化计算与业务含义联系起来。业务定义提供三条领域前提：（B1）选中的事实行具有非空键 $K$，其取值与业务事件一一对应；（B2）$e$ 求值得到每个事件的度量值；（B3）时间绑定与定义过滤表达预期的业务选择。数据检查在这些前提下建立计算所需的结构性质。}
\begin{proposition}[\bt{Sufficiency}{充分性}]
\label{prop:sufficiency}
\bi{If B1--B3 and every $c\in C(m^r)$ hold on $D_s$, then for any period $p$ each retained qualifying event appears exactly once in the input to $\mathrm{AGG}(e)$, and the unmatched fraction on each checked table or path is bounded by its publication baseline plus 0.1 percentage points. When all required date and inner joins match completely, the aggregate is exact for $p$.}
{若 B1--B3 及每个 $c\in C(m^r)$ 在 $D_s$ 上成立，则对任意期间 $p$，每个保留下来的合格事件在 $\mathrm{AGG}(e)$ 的输入中恰好出现一次，且每个受检查表或路径的未匹配比例不超过发布时基线加 0.1 个百分点。若全部所需日期连接和内连接都完整匹配，则该期间的聚合结果精确。}
\end{proposition}
\begin{proof}[\bt{Proof sketch}{证明概要}]
\bi{B1 and filtered key uniqueness give one fact row per event. Join cardinality prevents duplication; date and inner joins may remove unmatched rows. B2 and B3 assign each retained row its measurement and intended period. Complete matching retains all qualifying events, while completeness predicates bound the unmatched fraction on their respective tables and paths.\ifbilingual\else\qedhere\fi}
{B1 与过滤后键唯一性保证每个事件有一条事实行。连接基数防止行被复制；日期连接与内连接可能移除未匹配的行。B2、B3 将每条保留的行归入预期期间并提供其度量值。完整匹配保留全部合格事件；完整性谓词约束各受检查表和路径的未匹配比例。\qedhere}
\end{proof}
\bi{Completeness bounds apply to the checked populations, not separately to each reporting period. Table- and path-level checks can therefore be shared across period queries, while exactness for a period requires complete matching as stated above. The orange condition in Figure~\ref{fig:architecture}, $\forall c\in C(m^r):c(D_s)$, is the conjunction that the execution path must establish.}
{完整性界限作用于受检查的总体，而不是分别作用于每个报表期间。因此，表级和路径级检查可以跨期间查询共享；期间聚合的精确性则要求上述完整匹配。图~\ref{fig:architecture} 中橙色框内的 $\forall c\in C(m^r):c(D_s)$，正是执行路径需要建立的条件合取。}

\subsection{\bt{Lookup and shared maintenance}{查找与共享维护}}
\label{sec:maintenance}
\bi{The metric lookup interface, \code{find\_metric}, matches names or phrases against canonical names and aliases extracted during learning. It compares stored dependency versions and referenced join entries with their current state. Lookup distinguishes \emph{valid} when required evidence matches, \emph{pending} when revalidation is needed, and \emph{invalid} when a condition failed. Pending denotes missing current evidence. Available revisions return the definition, structured fields, generated SQL, and revision number. Unavailable entries retain business meaning and historical mappings as repair hints, without authorizing their execution. Agents issue at most three memory lookups per question before standard schema exploration.}
{指标查找接口 \code{find\_metric} 将名称或短语与学习时提取的规范名称、别名匹配，并比较所存依赖版本、引用连接与当前状态。所需证据匹配时为\emph{有效}，需要重验时为\emph{待验证}，条件失败时为\emph{失效}；待验证表示缺少当前证据。可用版本返回定义、结构化字段、生成的 SQL 和修订号。不可用条目仍保留业务含义和历史映射作为修复线索，不授予其执行权限。每题最多查找三次记忆，之后回到常规模式探索。}

\bi{The Version tracker supplies the table state used for dependency comparison. The lookup implementation can combine a schema fingerprint, DML counters, and an ETL batch identifier in $V(T)$. Snapshot-bound execution additionally uses transactional counters, defined in Section~\ref{sec:execution}, because asynchronous statistics do not identify the state visible to a read transaction. For a set of tables $\mathcal T$, $V(\mathcal T)$ denotes the vector of their versions in a fixed table order. A changed or missing shared join entry also makes its dependent definition pending, even when table versions alone match.}
{版本跟踪组件提供依赖比较所需的表状态。查找实现可以将结构指纹、DML 计数与 ETL 批次号组合为 $V(T)$。快照绑定执行进一步使用第~\ref{sec:execution} 节定义的事务性计数器，因为异步统计不能标识读事务所见的数据状态。对于表集合 $\mathcal T$，$V(\mathcal T)$ 表示按固定表顺序排列的版本向量。共享连接条目变化或消失时，其依赖定义同样变为待验证，即使单独比较表版本仍然匹配。}

\bi{\emph{Sharing condition evidence.} Let $\operatorname{id}(c)$ be the normalized identity of the executable check, including its type, tables, keys, join direction, filters, and null policy. The shared result store $\mathcal H$ indexes check statistics by}
{\emph{共享条件证据。}设 $\operatorname{id}(c)$ 为可执行检查的规范化标识，包含类型、表、键、连接方向、过滤和空值规则。共享结果库 $\mathcal H$ 按下式索引检查统计量：}
\begin{equation}
\label{eq:condition-result-key}
k(c,V)=\bigl(\operatorname{id}(c),V(\operatorname{dep}(c))\bigr).
\end{equation}
\bi{An outcome $o$ contains the statistics returned by $Q_c$ and the check's verdict. We use $\operatorname{Accept}(c,o)$ to apply the condition's acceptance rule; completeness conditions additionally compare the measured loss with their definition's recorded bound. Sharing the statistics therefore preserves each consumer's own bound. In Figure~\ref{fig:architecture}, the three references to $c_2$ find the same outcome at $V(T_3)$, whereas the write to $T_2$ changes the key for $c_3$.}
{结果 $o$ 包含 $Q_c$ 返回的统计量及检查判定。$\operatorname{Accept}(c,o)$ 表示应用该条件的接纳规则；完整性条件还将测得的丢行比例与所属定义记录的界限比较。因此，共享统计量仍保留各使用者自己的界限。图~\ref{fig:architecture} 中三个对 $c_2$ 的引用在 $V(T_3)$ 上找到同一结果，而对 $T_2$ 的写入改变了 $c_3$ 的结果键。}

\bi{Algorithm~\ref{alg:maintenance} implements steps 5--7 for a requested revision $x$, where $x$ is either a metric or a feature and $\bot$ denotes unavailability. Unchanged conditions retain their evidence. For an affected condition, Maintain first checks $\mathcal H$ at the required versions and runs $Q_c$ only on a miss. Concurrent requests for the same missing check join one in-flight execution; requests maintaining the same definition similarly join its current maintenance operation. The owner refreshes referenced joins and records the resulting dependency stamps before making the revision available. Thus, an update is handled by its distinct missing checks and affected definitions, rather than by repeating the work for every arriving agent.}
{算法~\ref{alg:maintenance} 对请求的版本 $x$ 实现第 5--7 步，$x$ 可以是指标或特征，$\bot$ 表示不可用。未变化条件保留原有证据。对于受影响条件，维护组件先在所需版本上查找 $\mathcal H$，仅未命中时才执行 $Q_c$。请求同一缺失检查的并发使用者合并到同次执行；维护同一定义的请求也合并到该定义当前的维护操作。负责请求刷新引用的连接，并记录更新后的依赖版本，再将版本置为可用。因此，一次更新按不同的缺失检查和受影响定义处理，无需对每个到达的智能体重复工作。}

\begin{algorithm}[t]
\caption{\bt{Lookup with shared maintenance}{带共享维护的查找}}
\label{alg:maintenance}
\normalsize
\begin{algorithmic}[1]
\Require Requested revision $x$; shared result store $\mathcal H$
\Ensure Available revision with refreshed evidence, or $\bot$
\State Await ongoing maintenance; reread $x$
\State \textbf{if} $x$ is missing or invalid \textbf{then return} $\bot$
\State $V\gets$ current dependency versions
\State \textbf{if} stamps and joins match \textbf{then return} $x$
\State Join or start maintenance of $x$
\State \textbf{if} not owner \textbf{then return} owner's result
\State Reread $x$; if superseded, return usable revision or $\bot$
\State Refresh affected shared joins
\For{each affected $c\in C(x)$, added assumptions last}
  \State \textbf{if} required join unavailable \textbf{then}
  \State \hspace{\algorithmicindent} Invalidate $x$; \Return $\operatorname{Recover}(x,c)$
  \State $k\gets\bigl(\operatorname{id}(c),V(\operatorname{dep}(c))\bigr)$
  \State $o\gets\mathcal H[k]$; on miss, run/join $Q_c$ at $V$ and cache $o$
  \If{$\neg\operatorname{Accept}(c,o)$}
    \State Invalidate $x$; \Return $\operatorname{Recover}(x,c)$
  \EndIf
\EndFor
\State Record refreshed stamps and joins; \Return $x$
\end{algorithmic}
\end{algorithm}


\bi{\emph{Repair and regression.} The algorithm's $\operatorname{Recover}$ operation first handles rollback of an optimized metric when only its added assumption fails (Section~\ref{sec:optimization}). For other metric failures, Metric repair searches a restricted space. If grain uniqueness fails, it considers equality filters on columns with few distinct values. A passing filter must restore one row per business key without losing keys, and exactly one filter must pass; competing primary and backup copies would otherwise be indistinguishable. The filter is conjoined with the definition's existing filters. A revalidated join can likewise contribute an additional required filter.}
{\emph{修复与回归。}算法的 $\operatorname{Recover}$ 操作首先处理仅新增前提失败时的优化指标回退（第~\ref{sec:optimization} 节）。对于其他指标失败，指标修复在受限空间中搜索。粒度唯一性失败时，系统考虑低基数列上的等值过滤。通过的过滤须恢复每个业务键一行且不丢键，并且恰好只有一个过滤通过；否则，主数据与备份副本等竞争候选无法区分。该过滤与定义已有过滤合取。重新验证的连接同样可以提供新增的必需过滤。}

\bi{The repair becomes a candidate $m^{r+1}$ rather than overwriting the failed computation. It must pass the static, grain, and read-only SQL checks, then reproduce the verified learning computation through an as-of query on the learning snapshot. The regression diamond in Figure~\ref{fig:architecture} represents this final test. A passing revision is published once and serves all consumers; an ambiguous or rejected repair leaves the metric invalid. Feature failures are revoked rather than sent through scalar-answer repair. Subsequent successful sessions can relearn unavailable definitions through the publication path.}
{修复生成候选 $m^{r+1}$，而不是覆盖失败的计算。候选须通过静态、粒度和只读 SQL 检查，再通过学习快照上的 as-of 查询重现已核实的学习计算。图~\ref{fig:architecture} 中回归菱形表示这一最终测试。通过的版本只发布一次，供全部使用者复用；歧义或被拒绝的修复使指标保持失效。特征条件失败时直接撤销，不进入基于标量答案的修复。后续成功会话可以经发布路径重新学到不可用的定义。}

\subsection{\bt{Snapshot-bound metric execution}{快照绑定的指标执行}}
\label{sec:execution}
\label{sec:snapshot}
\bi{A write may commit after lookup, so the orange execution path in Figure~\ref{fig:architecture} rechecks evidence. Given a revision and question parameters, \code{run\_metric} compiles the stored fields, qualifies filters by their owning table, and attaches the revision to the query. Custom SQL declares its revisions through \code{run\_sql}. Both paths check availability, revision references, grain filters, and read-only SQL safety, then validate conditions on the query's execution snapshot. Proposition~\ref{prop:sufficiency} applies to compiled aggregates and their period/ranking compositions; manual changes still undergo SQL review.}
{查找之后仍可能有写入提交，因此图~\ref{fig:architecture} 中的橙色执行路径重新检查证据。给定修订号和题目参数，\code{run\_metric} 编译所存字段，将过滤列绑定到所属表，并为查询附上修订引用。自定义 SQL 通过 \code{run\_sql} 声明所用修订。两条路径均检查可用性、修订引用、粒度过滤与只读 SQL 安全性，再在查询执行快照上验证条件。命题~\ref{prop:sufficiency} 适用于编译聚合及其期间、排名组合；手动改写仍接受 SQL 审查。}

\bi{\emph{Transactional versions.} Let $v_s(T)$ be the transactional version of $T$ visible in snapshot $s$. A statement-level trigger increments one metadata counter for each modified table and writing statement, in the same transaction as the write. The implementation distributes increments over 16 counter shards and takes their sum as the table version, reducing contention between writers. Commits and rollbacks therefore affect data and counters atomically. Reading the transactional version through the query's multi-version concurrency control (MVCC) snapshot binds its stamp to the visible table data.}
{\emph{事务性版本。}设 $v_s(T)$ 为快照 $s$ 中可见的 $T$ 的事务性版本。语句级触发器对每个被修改的表、每条写语句递增一个元数据计数器，该操作与数据写入处于同一事务。实现将递增分散到 16 个计数分片，以分片之和作为表版本，减少写入方之间的争用。因此，提交和回滚原子地作用于数据与计数器。通过查询的多版本并发控制（MVCC）快照读取事务性版本，使版本标记与可见的表数据绑定。}

\bi{Algorithm~\ref{alg:snapshot-execution} takes a read-only query $q$, its declared revisions $\mathcal R$, and the transactionally stamped result store $\mathcal H_{\mathrm{tx}}$. This namespace is separate from results tagged only with lookup versions. For fixed revision metadata during condition selection, it collects $\mathcal C_{\mathcal R}=\bigcup_{m^r\in\mathcal R}C(m^r)$ and opens a read-only repeatable-read transaction. Each condition looks up an outcome using the snapshot's vector $v_s(\operatorname{dep}(c))$; a miss executes $Q_c$ on $D_s$ and shares its outcome with requests at the same versions. Conditions sharing a query reuse its statistics, while their acceptance rules all remain required. The implementation can apply the strictest completeness bound when combining such conditions. Only after every condition passes does $q$ execute on $D_s$. A failure aborts the query and invalidates the cached version observation so the next request enters normal maintenance.}
{算法~\ref{alg:snapshot-execution} 接收只读查询 $q$、其声明的版本集合 $\mathcal R$，以及带事务版本标记的结果库 $\mathcal H_{\mathrm{tx}}$。这一命名空间与仅带查找版本标记的结果分开。在条件选择期间修订元数据保持固定的前提下，算法收集 $\mathcal C_{\mathcal R}=\bigcup_{m^r\in\mathcal R}C(m^r)$，开启只读可重复读事务。每个条件按快照的版本向量 $v_s(\operatorname{dep}(c))$ 查找结果；未命中时，在 $D_s$ 上执行 $Q_c$，并向相同版本上的请求共享结果。共用查询的条件复用统计量，但各自的接纳规则都必须满足；实现可在合并时采用最严格的完整性界限。全部条件通过后，才在 $D_s$ 上执行 $q$。失败则中止查询，并使缓存的版本观察失效，让下一请求进入常规维护。}

\begin{algorithm}[t]
\caption{\bt{Snapshot-bound metric execution}{快照绑定的指标执行}}
\label{alg:snapshot-execution}
\normalsize
\begin{algorithmic}[1]
\Require Read-only SQL $q$; metric revisions $\mathcal R$; result store $\mathcal H_{\mathrm{tx}}$
\Ensure Query result $y$, or rejection
\State \textbf{if} revision, required-filter, or SQL review fails \textbf{then reject}
\State $\mathcal C_{\mathcal R}\gets\bigcup_{m^r\in\mathcal R}C(m^r)$
\State Open read-only repeatable-read $s$ with state $D_s$
\State Read transactional versions $v_s$ inside $s$
\For{each $c\in\mathcal C_{\mathcal R}$}
  \State $k\gets\bigl(\operatorname{id}(c),v_s(\operatorname{dep}(c))\bigr)$
  \State $o\gets\mathcal H_{\mathrm{tx}}[k]$; on miss, run/join $Q_c(D_s)$ and cache $o$
  \If{$\neg\operatorname{Accept}(c,o)$}
    \State Abort $s$; invalidate version observation; \textbf{reject}
  \EndIf
\EndFor
\State $y\gets q(D_s)$; commit $s$; \Return $y$
\end{algorithmic}
\end{algorithm}


\bi{Outcomes from an earlier snapshot can be reused because the transactional counters identify whether their dependency data changed. The following lemma justifies the cross-snapshot lookup in the algorithm.}
{事务性计数器能够标识依赖数据是否变化，因此来自更早快照的检查结果仍可能复用。下面的引理为算法中的跨快照查找提供依据。}
\begin{lemma}[\bt{Check-result reuse}{检查结果复用}]
\label{lem:reuse}
\bi{For read-only snapshots under the write-coverage assumptions below, if $v_{s_1}(T)=v_{s_2}(T)$ for every $T\in\operatorname{dep}(c)$, then $c(D_{s_1})=c(D_{s_2})$.}
{对于满足下述写入覆盖前提的只读快照，若对每个 $T\in\operatorname{dep}(c)$ 都有 $v_{s_1}(T)=v_{s_2}(T)$，则 $c(D_{s_1})=c(D_{s_2})$。}
\end{lemma}
\begin{proof}
\bi{The visible committed-transaction sets are nested in snapshot acquisition order, say $s_1\subseteq s_2$. A transaction that writes $T$ increments $T$'s counter in the same transaction. Equal versions leave no transaction in $s_2\setminus s_1$ that writes $T$, so both snapshots see the same state of every table that $Q_c$ reads.\ifbilingual\else\qedhere\fi}
{按快照获取顺序，可见的已提交事务集合相互包含，不妨设 $s_1\subseteq s_2$。写 $T$ 的事务在同一事务内递增 $T$ 的计数。版本相等意味着 $s_2\setminus s_1$ 中没有写 $T$ 的事务，因此两个快照看到 $Q_c$ 所读各表的同一状态。\qedhere}
\end{proof}
\bi{Checks executed outside the consuming transaction enter the transactional-version cache only when their before/after dependency counters match and no SQL-result cache supplies their rows. Otherwise, execution obtains fresh evidence inside its own snapshot. A write committed after snapshot acquisition is invisible to both the checks and $q$, closing the validation-to-execution gap illustrated by the two orange $D_s$ boxes.}
{在消费事务之外执行的检查，只有在检查前后依赖计数器相同、且检查行并非来自 SQL 结果缓存时，才进入事务性版本缓存。否则，执行层在自身快照内取得新证据。快照获取后提交的写入对检查和 $q$ 都不可见，由此消除图中两个橙色 $D_s$ 框所表达的验证与执行之间的间隙。}

% PostgreSQL 18: sql-copy.html (triggers/FREEZE), sql-createtrigger.html
% (foreign-key actions), sql-truncate.html (transactional but not MVCC-safe).
\bi{\emph{Write-coverage assumptions.} Statement-level triggers track inserts, updates, and deletes, including \code{COPY FROM} and cascading foreign-key actions, on unpartitioned base tables without inheritance. Counters are never reset or reused while cached checks remain valid. Non-MVCC operations, including \code{TRUNCATE}, table rewrites, and \code{COPY FREEZE}, require quiescent readers. After DDL, trigger-bypassing writes, or restore/reset operations, the deployment invalidates affected check results before reuse.}
{\emph{写入覆盖前提。}语句级触发器跟踪无分区、无继承的基表上的插入、更新和删除，包括 \code{COPY FROM} 与外键级联操作。缓存检查有效期间不重置或复用计数器。非 MVCC 操作，包括 \code{TRUNCATE}、表重写和 \code{COPY FREEZE}，要求没有在途读事务。发生 DDL、绕过触发器的写入或恢复／重置操作后，部署须在复用前使受影响的检查结果失效。}

\subsection{\bt{Shared optimization and rollback}{共享优化与回退}}
\label{sec:optimization}
\bi{Once a definition is published, optimizing its structured implementation benefits every consumer of the generated SQL. An optimization pass proposes new metric revisions while preserving the representation, compilation, and maintenance protocol above. Two deterministic rules generate candidates: \emph{key range} replaces the period join with a range predicate on the fact's date key, obtaining bounds through scalar subqueries on the date dimension; \emph{left-join removal} drops a left join that supplies neither a referenced column nor a filter. If no rule candidate is adopted, an optional LLM advisor examines the definition, generated SQL, execution plan, and table sizes. Passing proposals with rule-supported structure may be published; other passing proposals remain suggestions for review.}
{定义发布后，优化其结构化实现可使所有生成 SQL 的使用者受益。优化过程提出新的指标版本，并保留上述表示、编译与维护协议。两条确定性规则生成候选：\emph{日期键范围}将期间连接替换为事实表日期键上的范围谓词，通过日期维度上的标量子查询取得边界；\emph{去掉左连接}消除既不提供被引用列、也不提供过滤的左连接。若未采纳规则候选，可选的 LLM 顾问检查定义、生成的 SQL、执行计划与表规模。通过验证且结构受规则支持的提议可发布，其他通过检查的提议保留为待审阅建议。}

\bi{Publication checks three aspects in addition to the static, grain, and SQL checks. \emph{Assumptions}: key range adds a date-key-contiguity condition, checked before adoption and maintained thereafter. \emph{Equivalence}: the candidate and current revision must agree for the learning question as of learning time and, within one repeatable-read snapshot, for all twelve months of the learning year, a quarter, a period difference, and a monthly ranking. \emph{Cost}: execution times are measured in alternating candidate/current order over at least eight periods in that same snapshot. The 95\% confidence interval for the mean paired difference, candidate minus current, must lie entirely below zero, and the mean saving must be at least 10\%.}
{除了静态、粒度和 SQL 检查，发布还检查三方面。\emph{前提}：日期键范围规则增加日期键连续条件，采纳前检查，之后持续维护。\emph{等价}：候选与当前版本须对学习题按学习时刻返回相同结果，并在一个可重复读快照内对学习年份全部 12 个月、一个季度、一次期间差值和一次按月排名返回相同结果。\emph{代价}：在同一快照中，对至少 8 个期间交替测量候选与当前版本的执行时间。以候选减当前版本得到配对差值，其均值的 95\% 置信区间须整体低于 0，且平均节省至少 10\%。}
\begin{proposition}[\bt{Rule equivalence}{规则等价}]
\label{prop:rewrite}
\bi{Let the date dimension's key be unique on database state $D$. (a) If date-key contiguity holds on $D$, i.e., the keys of every month form an integer interval beginning one above the previous month's last key, the key-range revision and its predecessor agree for every period of consecutive months. (b) If the right side of a removed left join is unique on its join key on $D$, removing that join does not change the result.}
{设日期维度的键在数据库状态 $D$ 上唯一。（a）若日期键连续性在 $D$ 上成立，即每个月的键构成一个整数区间，且从上个月最后一个键加一开始，则键区间版本与前一版本对每个由连续月份组成的期间返回相同结果。（b）若被消除的左连接右侧在连接键上唯一，则消除该连接不改变结果。}
\end{proposition}
\begin{proof}[\bt{Proof sketch}{证明概要}]
\bi{(a) Contiguity makes a period's date keys exactly the integer interval between its bounds. A fact key belongs to that interval if and only if it joins one date row in the period. (b) A left join with at most one match retains each left row once and supplies nothing that is read.\ifbilingual\else\qedhere\fi}
{（a）连续性使一个期间的日期键恰好构成边界之间的整数区间。事实键落在该区间内，当且仅当它连接到该期间的一行日期。（b）至多一个匹配的左连接使每个左侧行恰好保留一次，且不提供任何被读取的内容。\qedhere}
\end{proof}
\bi{The premise in (a) becomes a condition of the optimized revision; the premise in (b) belongs to the previous revision and is no longer needed by the new one. The rule proof establishes equivalence wherever the maintained assumptions hold, while the sampled tests check the implementation and its cost. Maintain checks an added optimization assumption last. If it alone fails, all conditions shared with the predecessor have current passing evidence: the system invalidates the optimized revision and makes its predecessor current again. This is the rollback branch of $\operatorname{Recover}$ in Algorithm~\ref{alg:maintenance}.}
{（a）的前提成为优化版本的条件；（b）的前提属于前一版本，新版本不再需要它。规则证明在受维护前提成立时建立等价性，抽样测试则检查实现及其代价。维护最后检查新增优化前提。若只有它失败，与前一版本共有的全部条件均有当前版本上通过的证据，系统便使优化版本失效，并让前一版本重新成为当前版本。这对应算法~\ref{alg:maintenance} 中 $\operatorname{Recover}$ 的回退分支。}

\subsection{\bt{Entity-level feature definitions}{实体级特征定义}}
\label{sec:features}
\bi{Predictive tasks reuse the same fact, measure, grain, joins, and filters, but aggregate relative to individual prediction instances rather than a reporting period. RelBench's forecasting tasks use seed rows containing an entity and a seed time~\cite{robinson2024relbench}. A feature revision $f^r$ extends $I$ with an entity mapping, an optional look-back window, and a fill value. The entity mapping records the target table and key, the column reached through a foreign key or verified join, and an optional anchor such as a seed post's owner. Thus, the past-spending definition in Figure~\ref{fig:architecture} can be instantiated for each consumer's seed table.}
{预测任务复用同样的事实表、度量、粒度、连接与过滤，但围绕单个预测样本聚合，而不是报表期间。RelBench 的预测任务使用包含实体和种子时间的种子行~\cite{robinson2024relbench}。特征版本 $f^r$ 在 $I$ 上增加实体映射、可选的回看窗口和填充值。实体映射记录目标表与键、经外键或已验证连接到达的列，以及可选锚列，例如种子帖子的作者。因此，图~\ref{fig:architecture} 中的历史消费定义可针对各使用者的种子表实例化。}

\bi{Let seed $i$ contain entity key $u_i$ and cutoff $t_i$. For a joined fact row $x\in J_I(D_s)$, let $\eta(x)$ be its mapped entity key and $t(x)$ its event time. Let $a(u_i)$ be the seed's anchor value, or $u_i$ itself when no anchor is used. With a look-back window of $w$ days, compilation aggregates rows satisfying}
{设种子 $i$ 包含实体键 $u_i$ 与截止时间 $t_i$。对于连接后的事实行 $x\in J_I(D_s)$，设 $\eta(x)$ 为其映射后的实体键，$t(x)$ 为事件时间。设 $a(u_i)$ 为种子的锚列值；没有锚列时，$a(u_i)=u_i$。回看窗口为 $w$ 天时，编译聚合满足下式的行：}
\begin{equation}
\label{eq:feature-window}
\eta(x)=a(u_i),\qquad t_i-w<t(x)\leq t_i.
\end{equation}
\bi{Without a window, the lower bound is omitted. The compiler emits one lateral aggregate per definition for each seed and applies the recorded fill value when its aggregate is null. The upper bound $t\leq t_i$ in Figure~\ref{fig:architecture} is inclusive, matching the implementation; rows after the cutoff are excluded. This seed-time boundary is distinct from the database snapshot $D_s$: the snapshot selects a consistent database state, while the temporal predicate selects the history available to a prediction instance within that state.}
{未指定窗口时省略下界。编译器针对每个种子，为每个定义生成一个 LATERAL 聚合，并在聚合为空时应用记录的填充值。图~\ref{fig:architecture} 中上界 $t\leq t_i$ 包含截止时刻，与实现一致；晚于截止时刻的行被排除。这一种子时间边界不同于数据库快照 $D_s$：快照选择一致的数据库状态，时间谓词则在该状态内选择预测样本可使用的历史。}

\bi{\emph{Two reuse granularities.} A task program records a successful session's complete SQL, submitted definitions, and validation score. It supplies an executable starting point for retraining the same task, preserving expressions such as self-joins. Individual feature definitions expose components that can be selected, composed, and refined across compatible tasks sharing an entity. \code{find\_program} returns the latest task program's executable SQL and SQL compiled from its valid definitions; \code{find\_features} retrieves and compiles definitions for an entity; \code{compose\_features} compiles a chosen subset. Consumers validate candidate SQL on the current task and audit temporal leakage before adoption.}
{\emph{两种复用粒度。}任务方案记录成功会话的完整 SQL、所提交的定义及验证分数，为同一任务的重训提供可执行起点，保留自连接等表达式。单条特征定义则提供可在共享实体的兼容任务间选择、组合和细化的组件。\code{find\_program} 返回最近任务方案的可执行 SQL，以及由其有效定义编译的 SQL；\code{find\_features} 检索并编译某实体上的定义；\code{compose\_features} 编译选定子集。使用者在当前任务上验证候选 SQL，并在采用前审计时间泄漏。}

\bi{\emph{Publication and maintenance.} After a session produces features that outperform the entity-table reference, an extractor turns its SQL into definition drafts. Four checks precede publication: valid fields and many-to-one joins; uniqueness of the fact's grain key; value-by-value reproduction of the submitted column on 2,000 seeds by compiled SQL; and unchanged values when recomputed on views that hide rows after the seed time. Passing drafts enter the same shared definition store. Data updates follow Algorithm~\ref{alg:maintenance}, reusing matching check results and rechecking affected conditions. A failed feature condition revokes that definition; later successful sessions replenish the library through the same column-reproduction and temporal checks.}
{\emph{发布与维护。}一次会话产生的特征优于实体表参照后，提取器将其 SQL 转为定义草稿。发布前执行四项检查：字段与多对一连接合法；事实表粒度键唯一；编译 SQL 在 2,000 个种子上逐值重现提交的列；在隐去种子时间之后的行的视图上重算，值保持不变。通过的草稿进入同一共享定义库。数据更新按算法~\ref{alg:maintenance} 处理，复用匹配的检查结果，重验受影响条件。失败的特征条件导致该定义被撤销；后续成功会话通过同样的逐列重现与时间检查补充共享库。}
